%% ma: L12-B14-0094
%% lop: 12
%% chuong: 5
%% bai: 14
%% dang: 12.14.28
%% loai: TN4
%% muc: TH
%% dap_an: "D"
%% nguon: "Ảnh giáo viên gửi 07/11/2026 (Tư Duy Mở - tuduymo.com, đề tự luyện Phương trình mặt phẳng, đề lập phương trình mặt phẳng), Câu 8"
%% kien_thuc: "Bài 14 (phương trình mặt phẳng tổng quát)"
%% trang_thai: cho-duyet
%% ly_do: "Dạng 12.14.28 mới đề xuất, chờ giáo viên duyệt."
%% ghi_chu: "Tài liệu gốc không kèm đáp án; đáp án do Claude giải và kiểm bằng sympy."
\begin{ex}
Trong không gian $Oxyz$, cho điểm $A(a;b;2a-b+1)$ với các tham số thực $a$ và $b$. Tập hợp điểm $A$ khi $a$, $b$ thay đổi là một mặt phẳng có phương trình tổng quát
\choice{$x-2y-z+1=0$}{$2x-y-z+3=0$}{$2x-y+1=0$}{\True $2x-y-z+1=0$}
\loigiai{$x=a$, $y=b$, $z=2a-b+1=2x-y+1$ nên $A$ luôn thuộc mặt phẳng $2x-y-z+1=0$; ngược lại mọi điểm của mặt phẳng đều có dạng trên. Chọn D.}
\end{ex}
